\section{Tauber 漸近展開と相転移}

早期検知確率は
\[
  p_e(\beta)=\Prob(T_{\mathrm{det}}<T_d)=1-\mathcal{L}_{T_d}(\beta)
\]
である．ここで $\mathcal{L}_{T_d}$ は $T_d$ の Laplace 変換を表す．
低投資域 $\beta\to0+$ での振る舞いは，潜伏時間分布の尾の重さを反映する
\cite{feller1971,bingham1989,resnick2007}．

\begin{theorem}[Tauber 展開の剰余項]
$T_d\sim\Pareto(t_m,\alpha)$，$\alpha\in(0,1)$，$\beta>0$ のとき，
\begin{equation}
  p_e(\beta)=t_m^\alpha\Gamma(1-\alpha)\beta^\alpha-C(\beta),
  \label{eq:tauber}
\end{equation}
ただし
\[
  C(\beta)=\alpha t_m^\alpha\int_0^{t_m}(1-e^{-\beta t})t^{-(\alpha+1)}\dd t
\]
であり，
\begin{equation}
  0\le C(\beta)\le \frac{\alpha t_m}{1-\alpha}\beta .
  \label{eq:tauber-bound}
\end{equation}
従って $R(\beta)=p_e(\beta)-t_m^\alpha\Gamma(1-\alpha)\beta^\alpha=O(\beta)$ で，
相対誤差は $O(\beta^{1-\alpha})\to0$ である．
\end{theorem}

\begin{proof}
独立性から
\[
  p_e(\beta)=\int_{t_m}^{\infty}(1-e^{-\beta t})
  \alpha t_m^\alpha t^{-(\alpha+1)}\dd t .
\]
積分区間 $[t_m,\infty)$ を $[0,\infty)-[0,t_m]$ に分解し，
$u=\beta t$ と置換すると主要項
$t_m^\alpha\Gamma(1-\alpha)\beta^\alpha$ を得る．
補正項については $1-e^{-x}\le x$ を用いて
\[
  C(\beta)\le \alpha t_m^\alpha\beta\int_0^{t_m}t^{-\alpha}\dd t
  =\frac{\alpha t_m}{1-\alpha}\beta .
\]
\end{proof}

\begin{lemma}[Pareto Laplace 変換の漸近]
$T_d\sim\Pareto(t_m,\alpha)$，$\beta\to0+$ において，
\[
\E[T_de^{-\beta T_d}]
\sim
\begin{cases}
\alpha t_m^\alpha\Gamma(1-\alpha)\beta^{\alpha-1}, & \alpha\in(0,1),\\
-t_m\log(\beta t_m), & \alpha=1,\\
\alpha t_m/(\alpha-1), & \alpha>1.
\end{cases}
\]
\end{lemma}

\begin{proof}
密度表示から
\[
  \E[T_de^{-\beta T_d}]
  =\alpha t_m^\alpha\int_{t_m}^{\infty}t^{-\alpha}e^{-\beta t}\dd t .
\]
$u=\beta t$ とおくと，$\alpha\in(0,1)$ では
\[
  \E[T_de^{-\beta T_d}]
  =\alpha t_m^\alpha\beta^{\alpha-1}
  \Gamma(1-\alpha,\beta t_m)
  \sim \alpha t_m^\alpha\Gamma(1-\alpha)\beta^{\alpha-1}
\]
である．$\alpha=1$ では
\[
  \E[T_de^{-\beta T_d}]
  =t_m\int_{\beta t_m}^{\infty}\frac{e^{-u}}{u}\dd u
  =t_m E_1(\beta t_m)
  \sim -t_m\log(\beta t_m),
\]
ただし $E_1(x)=-\gamma-\log x+O(x)$ を用いた．$\alpha>1$ では
$t^{-\alpha}$ が $[t_m,\infty)$ 上可積分であるため，優収束定理より
\[
  \E[T_de^{-\beta T_d}]
  \to\alpha t_m^\alpha\int_{t_m}^{\infty}t^{-\alpha}\dd t
  =\frac{\alpha t_m}{\alpha-1}.
\]
\end{proof}

\begin{theorem}[保持境界なし残余損失の相転移]
保持境界を外す極限 $n\Delta\to\infty$ で
\[
  Q_{\mathrm{nb}}(\beta)
  =\eta(d_1\tau_R+d_0)\mathcal{L}_{T_d}(\beta)
  +\eta d_1\kappa_s\E[T_de^{-\beta T_d}]
  +L(1-\eta)\mathcal{L}_{T_d}(\beta)
\]
とおく．これは \eqref{eq:qfull} で $\Hbeta\to0$，
$\Mbeta\to\E[T_de^{-\beta T_d}]$ とした残余損失である．
$\kappa_s>0$ のとき，$\beta\to0+$ で
\[
Q_{\mathrm{nb}}(\beta)\sim
\begin{cases}
\eta d_1\kappa_s\alpha t_m^\alpha\Gamma(1-\alpha)\beta^{\alpha-1}\to\infty,
& \alpha\in(0,1),\\
-\eta d_1\kappa_s t_m\log(\beta t_m)\to\infty,& \alpha=1,\\
\eta(d_1\tau_R+d_0)+L(1-\eta)
+\eta d_1\kappa_s\alpha t_m/(\alpha-1)<\infty,& \alpha>1.
\end{cases}
\]
\end{theorem}

\begin{remark}
$\alpha\le1$ かつ $\kappa_s>0$ では，保持境界を外すと検知失敗後の
スキャン停止時間項が低投資域で発散する．有限の保持期間は
$\Mbeta$ を切断モーメントに置き換えるため，重尾潜伏下では
バックアップ設計が検知投資の補完として不可欠である．
\end{remark}
